% Reviewed inputs: planning/cutoff-appendix-20261108/WIDE_CORRECTION-v1.md
% and WIDE_CORRECTION_ADDENDUM-v1.md; exact source pins are in the release receipt.
\section{Derivations for the wide cutoff correction}
\label{pca:wide-details}

\subsection{The strict band and its prime terms}
\label{pca:wide-band-proof}

Use the original \(r,e_0,e_1,J,z_0\) on each side.
The primal-zero removal at the dyad \(J\le N\zeta<2J\)
implies \(D/(EE_1(Nh)^2)>Nz_0/4\). Hence
\[
 \frac{E N\zeta}{U^2}
 <\frac{4D}{U^2}\frac{Ne_0}{E_1(Nh)^2}
 \le\frac{14}{E_1Nh},
\]
using \(e_0\mid h\) and \(2D\le7U^2\).
For a band label \(T_1<Nm<E/T_2\) in
\eqref{pca:ordered}, \(m\ne1\), so the full inner
M\"obius sum vanishes. In its omitted part use \(c=m/a\).
The inner sum in \eqref{pca:ordered} becomes
\[
 -\sum_{\substack{c\mid m\\Nc<Nm/T_1}}\mu(c).
\]
Here \(1<Nm/T_1<\Omega<14\). In the Eisenstein field the
primes above two and three lie in \(S\). The remaining
inert primes have norm at least \(25\), and the only
possible good prime norms below \(14\) are \(7,13\).
Two good prime factors already have norm at least \(49\).
Consequently this sum includes the unit and precisely the
primes specified in \eqref{pca:wide-factor}.
Substituting \(n=d/m\) turns \(Nm<E/T_2\) into \(Nn>T_2\).
The upper bound on \(\Omega\) also proves the two stated
conditions for vanishing of the prime term and the band.

If a \(p\)-subtraction is present, \(Nm>T_1Np\) and
\(Nn>T_2\) give \(E N\zeta>Np\,U^2\ge2D\).
The original primal profile for that individual cutoff is
\[
 W\!\left(\frac{E E_1(Nh)^2Nz_0\,Nb}{D}\right),
 \qquad Nb\ge1 .
\]
Its argument exceeds two, since \(E_1(Nh)^2\ge Ne_0\).
The individual cutoff has zero primal contribution even
if its dyad survived the original primal-zero removal.
Its Poisson components are still present in the exact
paired selection.

\subsection{The fixed-intersection scalar estimate}
\label{pca:wide-scalar-proof}

Fix the actual row, primal and dual labels, ordered
factorizations, original dyads, and both deletion allocations.
Let \(\mathcal L\) be the union of the primes of \(S\), of the
row, of all these labels, and of \(\mathcal P\).
Its radical has polynomial norm after the inherited
\(N\mathfrak k\le D^5\) clip and the original support bounds.
Choose the exact intersection of \(d\) with \(\mathcal L\)
as a squarefree ideal \(\delta\), subject to all required
memberships and deletions. The good \(e_1\) factors,
conjugate \(d\)-overlaps, and actual new \(d\)-residue
supports are fixed by these choices. Thus the original
row phase is fixed on the free part of \(d\), and an
original zero stays zero. The free part avoids the full
union \(\mathcal L\).

In the product of the two copies of \eqref{pca:wide-factor},
their real factors \(\mu(d)\) leave the original outer
\(\mu(d)\). The step-step term needs one ideal variable
and a step-correction term uses \(d=mn\).
For two corrections, use the four cells of primes in \(d\)
specified by their membership in the two factors \(m_0,m_1\).
Both \(m_i\)'s and their complements are products of the
appropriate cells. This is an exact squarefree factorization
with pairwise coprime cells, including units. Expand the
finite small-prime choices and assign the fixed primes of
\(\delta\) among the cells, with at most
\(4^{\omega(\delta)}\) assignments. The tests \(p\mid m_i\)
are thereby fixed. The cell ideals have form \(\delta_j q_j\),
and the sign is \(\mu(\delta)\prod_j\mu(q_j)\).

Select the largest free cell by its norm, with deterministic
index order on ties, and fix the other free cells. If its
norm is \(t\), then \(t\ge(Y')^{1/4}\). Its interval lies
on the scale \(L=Y'/\prod_{\rm other}Nq_j\), with
\(L\gg(Y')^{1/4}\) for every nonempty selection.
Its deleted set \(\mathcal T\) consists of \(\mathcal L\)
and the primes in the other free cells, and still has
polynomial radical norm. Every factorization inequality
is an interval test in \(t\), since each \(m_i,n_i\)
either includes the selected cell or is fixed.
The tie order changes only an endpoint.
The original primal and direct tests and the middle
cutoffs also give finitely many interval tests.
At fixed labels \(X'\) is proportional to \(1/E\), the
middle parameter \(M\) is proportional to \(E\), and the
integer norm equality is fixed. The original strict and
nonstrict endpoints remain in the scalar sum.

For the fixed field \(\mathbb Q(\omega)\), write
\(M_{\mathcal T}(x)=\sum_{Nc\le x,\,(c,\mathcal T)=1}\mu(c)\).
The fixed-field deleted M\"obius input used here is
\[
 |M_{\mathcal T}(x)|\ll_\kappa
 x\log^C(2D)\exp(-c_\kappa\sqrt{\log(2x)})
\]
for \(x\ge D^\kappa\) and polynomial radical norm.
The underlying fixed-field zero-free input is
Thorner--Zaman, \emph{Algebra \& Number Theory} 13 (2019),
1039--1068, Theorem~3.1, specialized to the trivial
extension of this one fixed field.
The classical truncated Perron shift for the reciprocal
zeta function gives
\(M_{\varnothing}(x)\ll x\exp(-c\sqrt{\log(2x)})\).
A possible fixed exceptional zero affects the fixed
starting threshold.
Artificial deletions are transferred exactly by
\[
 M_{\mathcal T}(x)
   =\sum_{\substack{Na\le x\\
                   \operatorname{supp}a\subseteq\mathcal T}}
                              M_{\varnothing}(x/Na).
\]
All powers of the deleted primes occur in this convolution.
Split at \(Na=\sqrt x\). In the first range use the
fixed-field bound. In the second use the elementary ideal
count and \(1/Na\le x^{-\eta/2}(Na)^{-1+\eta}\), with
\(\eta\asymp1/\sqrt{\log(2x)}\).
The finite reciprocal Euler product at \(1-\eta\) costs
a power of \(\log(2D)\). Indeed, if the deleted radical has
norm \(\ll D^{C_0}\), split its primes at
\(\log^2(2D)\). Below that threshold the ideal count bounds
the reciprocal-power sum by
\(O((\log(2D))^{2\eta}\log\log(3D))\).
Above it use the bound \(O(\log(2D))\) for the sum of their
log norms. For \(x\ge D^\kappa\), these estimates give the
claimed uniform logarithmic Euler factor and subpower
decay. The underlying analytic zeta function stays fixed.

After factoring out the fixed scalar and phase, the weight
is \(E^{-2}G_W(c_0/E)\overline{G_W(c_1/E)}\) for fixed
positive \(c_i\). Its supremum plus integrated absolute
derivative is \(O(\mathcal M_{24}(W)^2Y^{-2})\) on each
remaining interval. This follows from boundedness of the
radial transform and its first log derivative, and the
fixed dyadic range of \(E\). Partial summation gives
\[
 \ll_\kappa \mathcal M_{24}(W)^2Y^{-2}L\log^C(2D)
                      \exp(-c_\kappa\sqrt{\log(2Y')}).
\]
Here \(L\gg(Y')^{1/4}\) permits the deleted bound with a
fixed smaller exponent in its polynomial lower cutoff.
The remaining free cells sum harmonically by the factors
\(1/Nq_j\) in \(L\). Together with the fixed-prime assignments
this proves \eqref{pca:wide-scalar}. It applies to a scalar
with the row and all other ideal labels fixed.

\subsection{The complete partner sum and the equality prefix}
\label{pca:wide-prefix-proof}

For each \(p\in\mathcal P\), use the coefficient
\[
 -\mu(\zeta)\mu(d)
 \sum_{\substack{mn=d,\;p\mid m\\Nm>T_1Np,\;Nn>T_2}}1
\]
on the partner and sum its original ordered factorizations.
This coefficient is independent of the partner frequency
and divisor bounded. The individual primal zero above
proves \(V_d^{[p]}=0\). The exact original Poisson identity
therefore gives the full complement
\(-\mathcal E_d^{[p]}\), with the same middle projection.
Recombining all modes includes the re-expressed direct
part and preserves every original raw value.

For the elementary equality bound, at fixed \(r,e_0,e_1,J\)
put \(L_0=D/(E E_1(Nh)^2)\) and \(J'=J/Ne_0\).
The equal-norm ideals have \(z_0=g\ell\), \(\mathfrak k=g\bar\ell\).
There are \(O(J'\log(2+2J'))\) choices on the dyad, up to
the original constant-scale convention. Their norm \(Nz_0\)
appears inversely in the Poisson scalar, and the equal-norm
profile scale is comparable to \(L_0/R\).
Thus the pointwise row bound is at most a subpower times
\(\mathcal M_{24}(W)L_0 R^{-1/2}(1+L_0/R)^{-4}\).
The \(O(R)\) original rows give squared counting norm
\(O((DH)^\epsilon\mathcal M_{24}(W)^2R L_0)\).
Summing the square roots over the original \(r,e_0,e_1,J\)
uses the harmonic factors \(1/(Nr\,Ne_0)\) and fixed
divisor costs. After division by \(D\), the squared norm
is \(O((DH)^\epsilon\mathcal M_{24}(W)^2R/E)\).
The same proof handles the modified coefficient because it
is divisor bounded and row independent at each allocation.
Apply Cauchy to the original and modified middle projections
and sum harmonically in \(d\). This proves the target-size
full-frequency compensation of Section~\ref{pca:wide-core}.

For the pure comparable prefix, keep the correction on the
equality side. Its coefficient is divisor bounded on its
fixed labels. A partner correction, when included, is a
coefficient on \(x\) alone after fixing the other labels.
The insertion interval remains complete inside the original
eligibility. The sharp insertion proof with powerful freezing
and exact common-factor eligibility applies. Its relative
loss is \(D^{1/16}\) from the erased sieve factor
\(\Phi(A_1,B/Ns_1)/A_1\ll D^{1/8}\).
A nonzero equality correction implies \(E_0Nh_0<14\),
hence \(E_0(Nh_0)^2<196\). Therefore the equality radial
scale is \(y_0=D/(E E_0(Nh_0)^2 R)\gg D/(ER)\).
Use both half radial powers and keep the remaining
equality-side \(7/2\) decay. The fixed-\(d\) rectangle
is bounded by a constant times
\[
 (DH)^\epsilon\mathcal M_{24}(W)^2
 \frac{R}{E}\,D^{1/16}\left(1+\frac{D}{ER}\right)^{-7/2}.
\]
Count the exact original middle subtraction independently
at the same inherited \(R/E\) cost with the common factors.
For \(E\le D^{13/56}\), \(D/(ER)\gg D^{1/56}\) at the pure
scales. The remaining decay pays \(D^{1/16}\) because
\((7/2)(1/56)=1/16\). Harmonic summation gives the stated
prefix. The partner-only argument retains its own ratio
\(N\mathfrak k/Nz\), so the same automatic lower bound on
its radial scale is unavailable. The coherent masked wide
divisor sum remains open.
